\secciongrande{Mecanismos de ejecución de código en un navegador web la arquitectura del navegador}

\resumen{
  La interacción en un navegador comienza con la solicitud de un recurso y finaliza con su visualización. El proceso se adapta a prioridades como velocidad, fidelidad al contenido o seguridad. La mayoría de los navegadores comparten una arquitectura de referencia compuesta por componentes esenciales.
}

\definicion{Arquitectura de referencia}{Conjunto de componentes básicos comunes a la mayoría de los navegadores web, independientemente de su implementación específica.}
\subtituloA{Ciclo de interacción del usuario}

\lista{
  \textbf{Inicio:} El usuario indica la dirección del recurso (url o uri). |
  \textbf{Finalización:} Visualización exitosa del contenido solicitado. |
  \textbf{Dependencias:} El proceso se configura según el propósito (ej. rapidez, fidelidad, seguridad de la comunicación con el servidor).
}
